\documentclass[freefonts,nologo]{amnestyreport}

\usepackage{booktabs}
\usepackage{tabularx}
\usepackage{enumitem}
\usepackage{xcolor}

\title{Community Safety and Human Rights}
\subtitle{A practical evidence review for local decision-makers}
\author{Northbridge Research Unit}
\reportyear{2026}
\reportmonth{August}
\reportindex{NRU/CSHR/2026}
\blurb{An editable report starter with an executive summary, findings, recommendations, and an evidence register.}

\begin{document}
\frontmatter
\maketitle
\tableofcontents

\mainmatter
\section{Executive Summary}

This report examines how a local service partnership can make public spaces safer
while protecting the rights of the people who use them. The fictional Northbridge
partnership combines resident observations, service data, and structured interviews
into a small set of actions that can be measured over the next six months.

\begin{quote}
The central recommendation is to pair practical environmental improvements with
transparent review: better lighting and wayfinding should be delivered alongside
clear escalation routes, public reporting, and an independent complaints channel.
\end{quote}

\section*{Priority actions}
\begin{enumerate}[leftmargin=*,itemsep=4pt]
  \item Publish a monthly dashboard covering response times, unresolved reports,
        and the status of agreed improvements.
  \item Introduce a named contact for residents and community organisations,
        with a target response within five working days.
  \item Review all enforcement activity for consistency, proportionality, and
        possible unequal impact before expanding the programme.
\end{enumerate}

\section{Scope and Method}

The review covers three neighbourhood corridors and the evening period from
18:00 to 01:00. It uses a deliberately small evidence set so that the same
method can be repeated each month. Replace the fictional values below with your
own sources and include a short note about sampling, limitations, and consent.

\begin{table}[ht]
  \centering
  \caption{Evidence collected for the pilot review}
  \begin{tabularx}{0.9\textwidth}{@{}l r X@{}}
    \toprule
    Source & Count & Use in the review \\
    \midrule
    Resident interviews & 42 & Experiences, priorities, and suggested remedies \\
    Service records & 186 & Response times and recurring locations \\
    Evening observations & 18 & Lighting, visibility, and access conditions \\
    Community submissions & 27 & Issues not visible in administrative data \\
    \bottomrule
  \end{tabularx}
\end{table}

\section{Findings}
\section{Access and visibility}

Residents consistently described the eastern footpath as difficult to navigate
after dark. The issue was not simply the number of lights: blocked sight lines,
unclear entrances, and missing signs made it hard to know which route was open
and where help could be found.

\section{Communication and trust}

People were more willing to report problems when they knew what would happen next.
The existing process gave no public indication of ownership or expected timing,
so repeated reports often looked like separate failures rather than progress.

\section{Accountability}

The pilot should publish both successful and unsuccessful outcomes. A dashboard
that reports only completed actions will overstate impact and make it harder for
residents to identify unresolved risks.

\section{Recommendations}

\begin{description}[leftmargin=2.8cm,style=nextline]
  \item[Within 30 days] Name a programme owner, publish a contact route, and
    establish the baseline measures in the evidence register.
  \item[Within 90 days] Repair the three highest-priority lighting faults,
    install consistent signs, and publish the first monthly dashboard.
  \item[Within 180 days] Commission an independent review of outcomes and
    publish a response to every recommendation.
\end{description}

\section{Evidence Register}

\begin{tabularx}{\textwidth}{@{}l X X@{}}
  \toprule
  ID & Claim to verify & Owner / next check \\
  \midrule
  E1 & Reported incidents fell after the lighting repairs. & Operations; compare monthly series \\
  E2 & Residents received a response within five working days. & Partnership lead; sample 20 cases \\
  E3 & Complaints were reviewed for consistency and proportionality. & Independent reviewer; publish summary \\
  \bottomrule
\end{tabularx}

\section*{Using this starter}
Replace the sample narrative with your own findings, retain the method note, and
keep the evidence register close to the claims it supports. The \texttt{freefonts}
and \texttt{nologo} options keep this example independent of Amnesty branding and
restricted fonts.
\end{document}
